\begin{theorem} \label{theorem:limit_of_one_plus_one_over_n_to_the_n_exists}
    The sequence $(a_n)_{n=1}^\infty$ defined by $a_n := (1 + 1/n)^n$ \CrefAndHyperrefIfExist{definition:convergence_and_limit_of_sequence_in_extended_metric_space_valued_in_ordered_commutative_monoid_with_adjoined_infinity}{converges} in $\mathbb{R}$ to a finite limit. That is, the \CrefAndHyperrefIfExist{definition:euler_number_e_as_limit}{Euler's number}
    $$e = \lim_{n \to \infty} \left(1 + \frac{1}{n}\right)^n = e$$
    exists.
\end{theorem}

\begin{proof}
    By \Cref{lemma:sequence_one_plus_one_over_n_to_the_n_is_strictly_increasing}, the sequence $(a_n)$ is strictly increasing, hence in particular it is \CrefAndHyperrefIfExist{definition:monotone_sequence_of_elements_in_a_partially_ordered_set}{nondecreasing}. By \Cref{lemma:sequence_one_plus_one_over_n_to_the_n_is_bounded_above_by_3}, the sequence $(a_n)$ is bounded above by $3$. Since $(a_n)$ is a nondecreasing sequence in $\mathbb{R}$ that is bounded above, it converges to its supremum by the \CrefAndHyperrefIfExist{corollary:monotone_convergence_theorem_for_sequences_of_real_numbers}{monotone convergence theorem for sequences of real numbers}. The limit is therefore a finite real number.
\end{proof}
